
\documentclass{article}
\usepackage{colortbl}
\usepackage{makecell}
\usepackage{multirow}
\usepackage{supertabular}

\begin{document}

\newcounter{utterance}

\centering \large Interaction Transcript for game `cladder', experiment `full\_v1.5\_default', episode 2340 with llama{-}3.1{-}8b{-}instruct.
\vspace{24pt}

{ \footnotesize  \setcounter{utterance}{1}
\setlength{\tabcolsep}{0pt}
\begin{supertabular}{c@{$\;$}|p{.15\linewidth}@{}p{.15\linewidth}p{.15\linewidth}p{.15\linewidth}p{.15\linewidth}p{.15\linewidth}}
    \# & \multicolumn{2}{c}{Player} && \multicolumn{2}{c}{Game Master} \\
    \hline

    \theutterance \stepcounter{utterance}  
    & & & \multicolumn{4}{p{0.6\linewidth}}{
        \cellcolor[rgb]{0.9,0.9,0.9}{
            \makecell[{{p{\linewidth}}}]{
                \texttt{\tiny{[P1$\langle$GM]}}
                \texttt{You are an expert in causal inference. The following question is not a typical commonsense query, but rather a meticulously designed question created by a professor specializing in causal inference, intended to assess the students' mastery of the course content.} \\
\texttt{Imagine a self{-}contained, hypothetical world with only the following conditions, and without any unmentioned factors or causal relationships: Eating citrus has a direct effect on vitmain C. Vitmain C has a direct effect on scurvy. For patients not consuming citrus, the probability of scurvy is 73\%. For patients consuming citrus, the probability of scurvy is 54\%. Does eating citrus positively affect scurvy through vitmain C?} \\
\texttt{Start your answer with 'Yes' or 'No', followed by additional reasoning or evidence to support your explanation.} \\
            }
        }
    }
     \\ \\

    \theutterance \stepcounter{utterance}  
    & \multicolumn{4}{p{0.6\linewidth}}{
        \cellcolor[rgb]{0.561,0.737,0.561}{
            \makecell[{{p{\linewidth}}}]{
                \texttt{\tiny{[P1$\rangle$GM]}}
                \texttt{No.} \\
\\ 
\texttt{To assess whether eating citrus positively affects scurvy through vitamin C, we need to consider the potential causal relationships and their directions. } \\
\\ 
\texttt{First, let's establish the given causal relationships: } \\
\\ 
\texttt{1. Eating citrus → Vitamin C (direct effect)} \\
\texttt{2. Vitamin C → Scurvy (direct effect)} \\
\\ 
\texttt{We are also given the probabilities of scurvy for patients not consuming citrus (73\%) and those consuming citrus (54\%). } \\
\\ 
\texttt{To determine whether eating citrus positively affects scurvy through vitamin C, we need to consider the potential causal relationship between eating citrus and scurvy. } \\
\\ 
\texttt{One way to approach this is to use the backdoor criterion, which states that a path from the cause (eating citrus) to the effect (scurvy) is unconfounded if all backdoor paths (i.e., paths from the cause to the effect through intermediate variables) are blocked.} \\
\\ 
\texttt{In this case, the backdoor path from eating citrus to scurvy through vitamin C is not blocked, as we have a direct effect from vitamin C to scurvy. However, we need to consider the potential confounding effect of the probability of scurvy for patients not consuming citrus (73\%).} \\
\\ 
\texttt{To determine whether eating citrus positively affects scurvy through vitamin C, we can use the following formula:} \\
\\ 
\texttt{P(scurvy | citrus) = P(scurvy | vitamin C) * P(vitamin C | citrus) + P(scurvy | not citrus) * P(not vitamin C | citrus)} \\
\\ 
\texttt{where P(scurvy | citrus) is the probability of scurvy for patients consuming citrus (54\%), P(scurvy | vitamin C) is the probability of scurvy for patients with vitamin C (which we don't know), P(vitamin C | citrus) is the probability of vitamin C for patients consuming citrus (which we don't know), P(scurvy | not citrus) is the probability of scurvy for patients not consuming citrus (73\%), and P(not vitamin C | citrus) is the probability of not having vitamin C for patients consuming citrus (which we don't know).} \\
\\ 
\texttt{However, we can simplify the problem by using the following formula:} \\
\\ 
\texttt{P(scurvy | citrus) = P(scurvy | vitamin C) * P(vitamin C | citrus) + P(scurvy | not citrus) * P(not vitamin C | citrus)} \\
\\ 
\texttt{Since P(scurvy | citrus) = 0.54, P(scurvy | not citrus) = 0.73, and P(not vitamin C | citrus) = 1 {-} P(vitamin C | citrus), we can rewrite the formula as:} \\
\\ 
\texttt{0.54 = P(scurvy | vitamin C) * P(vitamin C | citrus) + 0.73 * (1 {-} P(vitamin C | citrus))} \\
\\ 
\texttt{To determine whether eating citrus positively affects scurvy through vitamin C, we need to consider the potential values of P(vitamin C | citrus) and P(scurvy | vitamin C).} \\
\\ 
\texttt{If P(vitamin C | citrus) = 1, then P(scurvy | vitamin C) = 0.54 / 1 = 0.54. However, this would imply that eating citrus has no effect on scurvy through vitamin C, as the probability of scurvy for patients consuming citrus would be the same as the probability of scurvy for patients with vitamin C.} \\
\\ 
\texttt{If P(vitamin C | citrus) < 1, then P(scurvy | vitamin C) would be greater than 0.54, as the probability of scurvy for patients consuming citrus would be a weighted average of the probability of scurvy for patients with vitamin C and the probability of scurvy for patients not consuming citrus.} \\
\\ 
\texttt{However, we are given that the probability of scurvy for patients consuming citrus (54\%) is less than the probability of scurvy for patients not consuming citrus (73\%). This suggests that eating citrus has a negative effect on scurvy, rather than a positive effect.} \\
\\ 
\texttt{Therefore, based on the given information, we can conclude that eating citrus does not positively affect scurvy through vitamin C.} \\
            }
        }
    }
    & & \\ \\

    \theutterance \stepcounter{utterance}  
    & & & \multicolumn{2}{p{0.3\linewidth}}{
        \cellcolor[rgb]{0.9,0.9,0.9}{
            \makecell[{{p{\linewidth}}}]{
                \texttt{\tiny{[GM$|$GM]}}
                \texttt{No} \\
            }
        }
    }
    & & \\ \\

    \theutterance \stepcounter{utterance}  
    & & & \multicolumn{2}{p{0.3\linewidth}}{
        \cellcolor[rgb]{0.9,0.9,0.9}{
            \makecell[{{p{\linewidth}}}]{
                \texttt{\tiny{[GM$|$GM]}}
                \texttt{no} \\
            }
        }
    }
    & & \\ \\

    \theutterance \stepcounter{utterance}  
    & & & \multicolumn{2}{p{0.3\linewidth}}{
        \cellcolor[rgb]{0.9,0.9,0.9}{
            \makecell[{{p{\linewidth}}}]{
                \texttt{\tiny{[GM$|$GM]}}
                \texttt{game\_result = WIN} \\
            }
        }
    }
    & & \\ \\

\end{supertabular}
}

\end{document}
